%Chargement des paquets
\usepackage{amsmath}
\usepackage{amsfonts}
\usepackage{amssymb}
\usepackage{enumerate}
\usepackage{mathtools}
\usepackage{bbm}
\usepackage{xparse, etoolbox}
\usepackage{enumerate}
\usepackage{enumitem}
\usepackage{mathabx}
\usepackage{babel}[french]

%environment
\usepackage{amsthm}
\usepackage{keytheorems}
\usepackage{hyperref} 

%Interface théorème
\newkeytheorem{lem}[
	name = Lemme
]

\newkeytheorem{prop}[
	name = Proposition
]

\newkeytheorem{thm}[
	name = Théorème
]

\newkeytheorem{coro}[
	name = Corollaire
]

\renewcommand*{\proofname}{Démonstration}

\theoremstyle{definition}
\newtheorem{defn}{Définition}

\theoremstyle{definition}
\newtheorem*{nota}{Notation}

\theoremstyle{definition}
\newtheorem*{limi}{Liminaire}

\theoremstyle{remark}
\newtheorem{rmq}{Remarque}

\theoremstyle{plain}
\newtheorem*{fait}{Fait}


%Ensembles de nombres
\newcommand{\N} {\mathbb{N}}
\newcommand{\Ne}{\N^\ast}
\newcommand{\Z} {\mathbb{Z}}
\newcommand{\Q} {\mathbb{Q}}
\newcommand{\R} {\mathbb{R}}
\newcommand{\Rb}{\overline{\mathbb{R}}}
\newcommand{\Rp}{\R_+}
\newcommand{\Rm}{\R_-}
\newcommand{\K} {\mathbb{K}}
\newcommand{\Ket} {\mathbb{K^{\ast}}}
\newcommand{\Cx}{\mathbb{C}}

%Ensembles, tribus
\newcommand{\A}{\mathbb{A}}
\renewcommand{\P}{\mathcal{P}}
\newcommand{\T}{\mathcal{T}}
\newcommand{\F}{\mathcal{F}}
\newcommand{\C} {\mathcal{C}}
\newcommand{\ouv}{\mathcal{O}}
\newcommand{\B}{\mathcal{B}}
\newcommand{\M}{\mathcal{M}}
\newcommand{\etag}{\mathcal{E}}
\newcommand{\etagOp}{\etag^{+}(\Omega)}
\newcommand{\etagO}{\etag(\Omega)}
%%Lp avant quotientage
\newcommand{\Lic}{\mathcal{L}^1}
\newcommand{\Lpc}{\mathcal{L}^p}
\newcommand{\Linfc}{\mathcal{L}^{\infty}}
\newcommand{\Lifull}{\Lic(\Omega, \B, m)}
\newcommand{\Lpfull}{\Lpc(\Omega, \B, m)}
\newcommand{\Linffull}{\Linfc(\Omega, \B, m)}
%%Lp après quotientage
\newcommand{\Li}{L^1}
\newcommand{\Ld}{L^2}
\newcommand{\Lp}{L^p}
\newcommand{\Linf}{L^{\infty}}
\newcommand{\Listd}{\Li(\Omega, m)} 
\newcommand{\Ldstd}{\Ld(\Omega, m)} 
\newcommand{\Lpstd}{\Lp(\Omega, m)} 

%Quelques opérateurs
\DeclareMathOperator{\codim}{codim}
\DeclareMathOperator{\rg}{rg}
\DeclareMathOperator{\tr}{tr}
\DeclareMathOperator{\vect}{Vect}
\DeclareMathOperator{\ima}{Im}
\DeclareMathOperator{\id}{Id}
\DeclareMathOperator{\sign}{sign}
\DeclareMathOperator{\ext}{ext}
\newcommand{\ind}{\mathbbm{1}}
\newcommand*{\spr}[2]{\left(\,#1\,\middle|\,#2\,\right)}
\newcommand{\interior}[1]{%
  {\kern0pt#1}^{\mathrm{o}}%
}
\DeclareMathOperator{\card}{card}
\DeclareMathOperator{\ppcm}{ppcm}

%Commandes de pure flemme
\newcommand{\veps}{\varepsilon}
\newcommand{\intab}{\int_{a}^{b}}
\newcommand{\intR}{\int_{-\infty}^{\infty}}
\newcommand{\intinf}{\int_{- \infty}^{+ \infty}}
\newcommand{\equi}{\Leftrightarrow}
\newcommand{\Kev}{$\K$-espace vectoriel }
\newcommand{\Kevs}{$\K$-espaces vectoriels }
\newcommand{\Rev}{$\R$-espace vectoriel }
\newcommand{\Revs}{$\R$-espaces vectoriels }
\newcommand{\Cev}{$\Cx$-espace vectoriel }
\newcommand{\Cevs}{$\Cx$-espaces vectoriels }
\newcommand{\evn}{(E, \norm{.})}
\newcommand{\interf}[2]{[#1,#2]}
\newcommand{\limb}{\lim\limits_}
\newcommand{\lebn}{\lambda_n}
\newcommand{\pp}{\text{~presque partout}}
\newcommand{\derivp}[2]{\frac{\partial #1}{\partial #2}}
\newcommand{\famiu}{(u_i)_{i \in I}}
\newcommand{\sev}{\text{~s.e.v.}}
\newcommand{\Mnp}{M_{n,p}(\K)}
\newcommand{\Mn}{M_{n}(\K)}
\newcommand{\tq}{\text{~t.q.~}}
\newcommand{\mq}{~montrer~que~}
\newcommand{\et}{\quad \text{et} \quad}
\newcommand{\ou}{\text{~ou~}}
\newcommand{\Kx}{\K[X]}


%Commandes de gars chelou
\renewcommand{\subset}{\subseteq}

%Commande restriction, ty duckduckgo
\def\restr#1#2{\mathchoice
              {\setbox1\hbox{${\displaystyle #1}_{\scriptstyle #2}$}
              \restraux{#1}{#2}}
              {\setbox1\hbox{${\textstyle #1}_{\scriptstyle #2}$}
              \restraux{#1}{#2}}
              {\setbox1\hbox{${\scriptstyle #1}_{\scriptscriptstyle #2}$}
              \restraux{#1}{#2}}
              {\setbox1\hbox{${\scriptscriptstyle #1}_{\scriptscriptstyle #2}$}
              \restraux{#1}{#2}}}
\def\restraux#1#2{{#1\,\smash{\vrule height .8\ht1 depth .85\dp1}}_{\,#2}} 

%Quelques normes
\DeclarePairedDelimiter\abs{\lvert}{\rvert}
\DeclarePairedDelimiter\labs{\bigl\lvert}{\bigr\rvert}
\DeclarePairedDelimiter\norm{\lVert}{\rVert}
\DeclarePairedDelimiterXPP\normi[1]{}\lVert\rVert{_1}{\ifblank{#1}{\:\cdot\:}{#1}}
\DeclarePairedDelimiterXPP\normd[1]{}\lVert\rVert{_2}{\ifblank{#1}{\:\cdot\:}{#1}}
\DeclarePairedDelimiterXPP\normp[1]{}\lVert\rVert{_p}{\ifblank{#1}{\:\cdot\:}{#1}}
\DeclarePairedDelimiterXPP\normq[1]{}\lVert\rVert{_q}{\ifblank{#1}{\:\cdot\:}{#1}}
\DeclarePairedDelimiterXPP\norminf[1]{}\lVert\rVert{_\infty}{\ifblank{#1}{\:\cdot\:}{#1}}

%Bracket en angle
\DeclarePairedDelimiter\angl{\langle}{\rangle}

%Set, parenthèses
\newcommand{\set}[1]{\{ #1 \}}
\newcommand{\bigset}[1]{\bigl\{ #1 \bigr\}}
\newcommand{\Bigset}[1]{\Bigl\{ #1 \Bigr\}}
\newcommand{\bigpar}[1]{\bigl( #1 \bigr)}
\newcommand{\Bigpar}[1]{\Bigl( #1 \Bigr)}

%Opitmisation des opérateurs de comparaison
\DeclarePairedDelimiter\tio{o (}{)}
\DeclarePairedDelimiter\gro{O (}{)}

\newcommand{\Znz}{\Z/n\Z}
\newcommand{\Znzx}{\Z/n\Z^{\times}}

\pagestyle{fancy}

\fancyhead[C]{}
\fancyhead[R]{}

\begin{document}

\underline{Source :} Rombaldi - Algèbre - page 285

\underline{Leçons :}
\begin{itemize}
\item Leçon 120 : Anneaux Z/nZ. Applications. 
\end{itemize}

\begin{lem}
\label{lem:1}
Soit $(n_j)_{1 \leq j \leq r}$ une suite de $r \geq 2$ entiers naturels distincts de 0 et 1.
\begin{enumerate}
\item si les $n_j$ sont deux à deux premiers entre eux, alors leur $\ppcm$ est leur produit ;
\item sinon, on a $\ppcm(n_1, \dots, n_r) < \prod_j n_j$.
\end{enumerate}
\end{lem}

\begin{proof}
\begin{enumerate}
\item Dans ce cas, il n'y a aucun facteur commun dans les décompositions en facteurs premiers des $n_j$,
or un multiple de $n = \prod_i p_i^{m_i}$ est de la forme $k \prod_i p_i^{m_i}$. On en déduit le résultat.
\item Dans ce cas, il existe $p_1$ premier apparaissant dans deux décompositions en facteurs premiers des $n_j$.
En notant $\prod_j n_j=\prod_i p_i^{m_i}$, on a $m_1 \geq 2$. 
Donc $\prod_i p_i^{m_i}$ avec $m_1=1$ est un multiple commun de tous les $n_j$ et il divise $\prod_j n_j$ d'où le résultat.
\end{enumerate}
\end{proof}

\begin{thm}
Soit $(n_j)_{1 \leq j \leq r}$ une suite de $r \geq 2$ entiers naturels distincts de 0 et 1 et soit $n = \prod n_j$.
Les entiers $(n_j)_{1 \leq j \leq r}$ sont deux à deux premiers entre eux si, et seulement si, les anneaux $\Znz$ et 
$\prod_j \Z/n_j\Z$ sont isomorphes.
Dans ce cas, l'application :
\[ \begin{array}{c c c c}
\phi : & \Znz & \to & \prod_{j=1}^r \Z/n_j\Z \\
& \pi_n(k) & \mapsto & (\pi_{n_1}(k), \dots, \pi_{n_r}(k))
\end{array} \]
est un isomorphisme d'anneau d'inverse :
\[ \begin{array}{c c c c}
\phi^{-1} : &  \prod_{j=1}^r \Z/n_j & \to & \Znz \\
& (\pi_{n_1}(a_1), \dots, \pi_{n_r}(a_r)) & \mapsto & \pi_n \left ( \sum_{j=1}^r a_j u_j \dfrac{n}{n_j} \right ) ,
\end{array} \]
où $(u_j)_{1 \leq j \leq r}$ est une suite d'entiers relatifs telle que $\sum_{j=1}^r a_j u_j \dfrac{n}{n_j} = 1$.
\end{thm}

\begin{proof}
Le produit cartésien $\prod_j \Z/n_j\Z$ est muni de la structure d'anneau produit.
\begin{enumerate}[label=(\roman*)]

\item
Supposons les $(n_j)_{1 \leq j \leq r}$ deux à deux premiers entre eux.
L'application :
\[ \begin{array}{c c c c}
f : & \Z & \to & \prod_{j=1}^r \Z/n_j\Z \\
& k & \mapsto & (\pi_{n_1}(k), \dots, \pi_{n_r}(k))
\end{array} \]
est un morphisme d'anneaux et son noyau est formé des entiers multiples de tous les $n_j$, donc de leur $\ppcm$ qui est égal à
$n=\prod_j n_j$ selon le lemme $\ref{lem:1}$. On peut donc le factoriser en un morphisme injectif :
\[ \begin{array}{c c c c}
\phi : & \Z/n\Z & \to & \prod_{j=1}^r \Z/n_j\Z \\
& \pi_n(k) & \mapsto & (\pi_{n_1}(k), \dots, \pi_{n_r}(k))
\end{array} \]
Par dénombrement direct, on remarque que ces deux anneaux ont le même cardinal, donc $\phi$ est un isomorphisme.

\item
Si les $(n_j)_{1 \leq j \leq r}$ ne sont pas deux à deux premiers entre eux, alors $d = \ppcm(n_1, \dots, n_r) < n$.
Alors, $\phi(\overline{d}) = (1, \dots, 1)$ et donc $\phi$ n'est pas injective, en particulier, les anneaux ne sont pas isomorphes.

\item
On suppose à nouveau que les entiers $(n_j)_{1 \leq j \leq r}$ sont deux à deux premiers entre eux.
On désigne par $(m_j)_{1 \leq j \leq r}$ la suite d'entiers définis par :
\[ m_j = \dfrac{n}{n_j} = \prod_{\substack{i=1 \\ i \neq j}}^r n_i , \]
de façon à ce que chaque $m_j$ soit divisible par tous les $n_i$ excepté $n_j$.
Les $m_j$ sont premiers entre eux dans leur ensemble. 
En effet, si $p$ premier divise tous les $m_j$ alors, il apparait dans les décompositions en facteurs premiers de tous les $n_j$ ce qui est
impossible.
On déduit, d'après l'identité de Bézout qu'il existe une suite $(u_i)_{1 \leq j \leq r}$ d'entiers tels que :
\[ \sum_{j=1}^r m_j u_j = 1 , \]
et par suite :
\[ \pi_j(1) = \pi_j \left ( \sum_{j=1}^r m_j u_j  \right ) = \pi_j(m_j) \pi_j(u_j) = 1 , \]
autrement dit, $ \pi_j(m_j)$ est inversible dans $\Z/n_j\Z$ et d'inverse $\pi_j(u_j)$.
Soit $(a_1, \dots, a_r) \in \prod_j \Z/n_j\Z$ on pose :
\[ k = \sum_{j=1}^r a_j u_j m_j , \]
on a :
\[ \pi_j(k) = \pi_j(a_j)\pi_j(u_j)\pi_j(m_j) = \pi_j(a_j) , \]
et donc :
\[ \phi(\pi_n(k)) = (\pi_1(k), \dots, \pi_r(k)) = (\pi_1(a_1), \dots, \pi_r(a_r)) . \]
On a donc bien démontré que l'isomorphisme inverse est donné par :
\[ \phi^{-1}(\pi_1(a_1), \dots, \pi_r(a_r)) = \pi_n(k) , \]
avec $k$ ayant la forme énoncé.
\end{enumerate}
\end{proof}

\begin{lem}
\label{lem:2}
Si $A, B$ sont deux anneaux unitaires et $\psi$ est un isomorphisme d'anneaux de $A$ sur $B$, 
alors $\psi$ réalise un isomorphisme de groupe de $(A^{\times}, \cdot)$ sur $(B^{\times}, \cdot)$.
\end{lem}

\begin{proof}
Si $a \in A^{\times}$ alors $\psi(aa^{-1}) = \psi(a) \psi(a^{-1}) = 1_B$ donc $\psi(a) \in A^{\times}$ 
et $\psi(a)^{-1} = \psi(a^{-1})$. Ainsi $\psi$ applique $A^{\times}$ dans $B^{\times}$.
Comme $\psi$ est injective sur $A$ elle l'est sur $A^{\times}$.
Soit $b \in B^{\times}$, il existe $c \in B^{\times}$ tel que $bc=1_B$. 
Par surjectivité de $\psi$, il existe $a, a' \in A$ tels que $b=\psi(a)$ et $c=\psi(a')$ et donc,
$\psi(a)\psi(a') = \psi(aa') = 1_B$ soit $aa' = 1_A$ par injectivité et $a \in A^{\times}$.
On a bien le résultat attendu.
\end{proof}

\begin{coro}
Soit $n \geq 2$ et $\prod_i p_i^{m_i}$ sa décomposition en facteurs premiers.
En notant $\varphi$ la fonction indicatrice d'euler, on déduit que :
\[ \varphi(n) = \prod_i \varphi(p_i^{m_i}) . \]
\end{coro}

\begin{lem}
Pour tout nombre premier $p$ et tout entier naturel $m$, on a $\varphi(p^m) = (p-1) p^{m-1}$.
\end{lem}

\begin{proof}
Parmi les nombres de 1 à $p^m$, seuls les multiples de $p$ ne sont pas premiers avec $p^m$, soit $p^{m-1}$ possibilités.
\end{proof}

\begin{thm}
Soit $n \geq 2$ et $\prod_i p_i^{m_i}$ sa décomposition en facteurs premiers.
On a alors :
\[ \varphi(n) = \prod_i (p_i-1) p_i^{m_i-1} = n \prod_i \left ( 1 - \dfrac1{p_i} \right ) . \]
\end{thm}

\end{document}